\documentclass[a4paper,11pt]{article}
\usepackage[utf8]{inputenc}
\usepackage[T1]{fontenc}
\usepackage[french]{babel}
\usepackage[top=2cm, bottom=2cm, left=2.2cm, right=2.2cm]{geometry}
\usepackage{xcolor}
\usepackage{graphicx}
\usepackage{booktabs}
\usepackage{tabularx}
\usepackage{enumitem}
\usepackage{listings}
\usepackage{float}
\usepackage{caption}
\usepackage{hyperref}
\usepackage{amssymb}
\usepackage{pifont}
\usepackage{tcolorbox}

\hypersetup{colorlinks=true, linkcolor=primary, urlcolor=primary}

\definecolor{primary}{HTML}{1A237E}
\definecolor{success}{HTML}{2E7D32}
\definecolor{warning}{HTML}{E65100}
\definecolor{danger}{HTML}{B71C1C}
\definecolor{lightgreen}{HTML}{E8F5E9}
\definecolor{lightorange}{HTML}{FFF3E0}
\definecolor{codebg}{HTML}{F5F5F5}
\definecolor{codeblue}{rgb}{0.13,0.29,0.53}
\definecolor{codegreen}{rgb}{0.0,0.50,0.0}
\definecolor{codegray}{rgb}{0.5,0.5,0.5}

\lstdefinestyle{sh}{
  backgroundcolor=\color{codebg},
  basicstyle=\ttfamily\footnotesize,
  keywordstyle=\color{codeblue}\bfseries,
  commentstyle=\color{codegreen}\itshape,
  numberstyle=\tiny\color{codegray},
  breaklines=true,
  frame=single,
  showstringspaces=false,
  tabsize=2,
  columns=fullflexible,
}
\lstset{style=sh}

\newcommand{\ok}{\textcolor{success}{\checkmark}}
\newcommand{\ko}{\textcolor{danger}{\ding{55}}}

\title{\textbf{Mise en service du LiDAR et visualisation RViz\\ (ProtoVA2)}\\[2mm]
\large Journal technique : diagnostic, panne matérielle, réparation et validation}
\author{Douae Sebti}
\date{19 juin 2026}

\begin{document}
\maketitle

\begin{abstract}
\noindent Ce document retrace, étape par étape, la mise en service du LiDAR du
ProtoVA2 (Jetson Orin Nano, ROS\,2 Humble) en vue de la conduite autonome.
Il inclut volontairement \textbf{toutes les tentatives, y compris celles qui ont
échoué}, car le cheminement du diagnostic constitue l'essentiel du travail : un
LiDAR qui « tournait » mais ne renvoyait aucune donnée. On y traite
l'identification du capteur, l'élimination méthodique des causes logicielles, la
découverte d'une panne matérielle (connecteur de données), sa réparation, puis
la validation du flux \texttt{/scan} et sa visualisation dans RViz sur le PC via
WiFi.
\end{abstract}

\tableofcontents
\bigskip

% =====================================================================
\section{Contexte et objectif}
Après la démonstration de téléopération sur WiFi (DDS natif), l'objectif était
de démarrer le volet \textbf{autonomie réactive}. Le premier capteur
indispensable est le \textbf{LiDAR} (télémètre laser rotatif) qui publie le
topic \texttt{/scan} : c'est la perception sur laquelle repose l'évitement
d'obstacles. Tout le travail ci-dessous a été réalisé à distance, par SSH depuis
le PC vers la Jetson (\texttt{protova2@10.37.11.28}).

% =====================================================================
\section{Identification du matériel}
Inspection des périphériques connectés à la Jetson :
\begin{lstlisting}[language=bash]
lsusb                       # liste des peripheriques USB
ls -l /dev/ttyUSB* /dev/ttyACM*
ls -l /dev/serial/by-id/
\end{lstlisting}

\noindent Résultats :
\begin{itemize}[leftmargin=1.4em]
  \item \texttt{ID 10c4:ea60 Silicon Labs CP210x UART Bridge} sur
        \textbf{\texttt{/dev/ttyUSB0}} : c'est le pont USB-série du LiDAR.
  \item \texttt{Raspberry Pi Pico} sur \texttt{/dev/ttyACM0} : carte de commande
        des actionneurs (à ne pas confondre).
  \item Caméra Orbbec Femto Bolt (USB) — sans rapport ici.
\end{itemize}
Le pilote \texttt{rplidar\_ros} est déjà compilé dans l'espace de travail
\texttt{protova2\_yass}. Le capteur est donc un \textbf{RPLIDAR} relié par
CP2102 ; il restait à déterminer le modèle (qui fixe le débit série).

% =====================================================================
\section{Le symptôme : aucune donnée}
Premier essai du pilote ROS au débit 115200 (RPLIDAR A1/A2) :
\begin{lstlisting}[language=bash]
ros2 run rplidar_ros rplidar_composition --ros-args \
  -p serial_port:=/dev/ttyUSB0 -p serial_baudrate:=115200 \
  -p frame_id:=laser_frame
\end{lstlisting}
\begin{lstlisting}
[INFO]  RPLIDAR SDK Version: 2.1.0
[ERROR] Error, operation time out. SL_RESULT_OPERATION_TIMEOUT!
[ros2run]: Process exited with failure 255
\end{lstlisting}
Le port s'ouvre, mais le capteur ne répond pas à la requête d'identification
(\texttt{getDeviceInfo}). \textbf{La tête du LiDAR tournait pourtant bien.}

% =====================================================================
\section{Diagnostic : élimination des causes (tentatives échouées)}
L'erreur \texttt{TIMEOUT} pouvant venir de multiples sources, chacune a été
testée puis écartée.

\subsection{(\ko) Mauvais débit série}
Sur une liaison série, les deux extrémités doivent parler à la même vitesse.
Les quatre débits standard des RPLIDAR ont été essayés :
\begin{center}
\begin{tabular}{@{}lll@{}}
\toprule
Débit & Modèle visé & Résultat \\
\midrule
115\,200 & A1 / A2     & \ko\ timeout \\
256\,000 & A3 / S1     & \ko\ timeout \\
460\,800 & C1          & \ko\ timeout \\
1\,000\,000 & S2 / S3  & \ko\ timeout \\
\bottomrule
\end{tabular}
\end{center}
\textbf{Conclusion :} échec identique à toutes les vitesses $\Rightarrow$ ce
n'est pas un problème de débit.

\subsection{(\ko) Droits d'accès au port}
\begin{lstlisting}[language=bash]
id protova2 | grep dialout      # -> 20(dialout)
\end{lstlisting}
L'utilisateur appartient au groupe \texttt{dialout} ; il a donc le droit de
lire/écrire sur le port (sinon l'erreur serait « accès refusé », pas un
timeout). \textbf{Cause écartée.}

\subsection{(\ko) Interférence de ModemManager}
\texttt{ModemManager} (service de gestion des modems, utilisé pour la 5G) sonde
parfois les ports série et perturbe la communication. Vérification :
\begin{lstlisting}[language=bash]
systemctl is-active ModemManager   # -> inactive
\end{lstlisting}
Le service est \textbf{inactif}. \textbf{Cause écartée.}

\subsection{(\ko) Sonde série bas niveau}
Pour savoir si \emph{quoi que ce soit} revient du capteur, un petit script
Python (\texttt{pyserial}) envoie directement les commandes du protocole RPLIDAR
et affiche la réponse brute. Une réponse valide commence par les octets
\texttt{A5 5A}.
\begin{lstlisting}[language=Python]
import time, serial
GET_INFO   = bytes([0xA5, 0x50])
GET_HEALTH = bytes([0xA5, 0x52])
STOP       = bytes([0xA5, 0x25])
for baud in (115200, 256000, 460800, 1000000):
    ser = serial.Serial('/dev/ttyUSB0', baud, timeout=1)
    ser.dtr = False; time.sleep(0.2); ser.reset_input_buffer()
    ser.write(STOP); time.sleep(0.05); ser.reset_input_buffer()
    ser.write(GET_INFO); time.sleep(0.3)
    print(baud, ser.read(64).hex(' ') or '(rien)')
    ser.close()
\end{lstlisting}
\begin{lstlisting}
[115200]  GET_INFO -> (rien)
[256000]  GET_INFO -> (rien)
[460800]  GET_INFO -> (rien)
[1000000] GET_INFO -> (rien)
\end{lstlisting}
\textbf{Aucun octet ne revient}, à aucun débit, même en contournant ROS.

\subsection{(\ko) Lignes de contrôle DTR / RTS}
Sur certains adaptateurs, le transceiver de données dépend de l'état des lignes
DTR/RTS. Les quatre combinaisons ont été testées (avec impulsion de reset) :
\begin{lstlisting}
DTR=0 RTS=0 -> (rien)   DTR=0 RTS=1 -> (rien)
DTR=1 RTS=0 -> (rien)   DTR=1 RTS=1 -> (rien)
\end{lstlisting}
\textbf{Cause écartée.}

% =====================================================================
\section{Conclusion du diagnostic : panne matérielle}
\begin{table}[H]
\centering
\begin{tabularx}{\textwidth}{@{}l c X@{}}
\toprule
Élément testé & État & Interprétation \\
\midrule
Énumération USB / ouverture du port & \ok & câble USB (adaptateur$\to$Jetson) bon \\
Permissions (\texttt{dialout})       & \ok & pas un problème de droits \\
ModemManager                         & \ok & inactif, pas d'interférence \\
4 débits série                       & \ok & pas une mauvaise config \\
Lignes DTR/RTS (4 combinaisons)      & \ok & pas un souci de ligne de contrôle \\
\textbf{Réponse données du LiDAR}    & \ko & \textbf{le capteur n'émet rien} \\
Rotation de la tête (moteur)         & \ok & \textbf{alimentation présente} \\
\bottomrule
\end{tabularx}
\end{table}

\noindent \textbf{Verdict :} la tête tourne (donc le LiDAR est \emph{alimenté})
mais ne renvoie \emph{aucune donnée}. C'est le signe que \textbf{les broches
d'alimentation font contact} (le moteur tourne) \textbf{mais pas les broches de
données}. Autrement dit, un \textbf{défaut de câble/connecteur de données}, pas
un problème logiciel ni d'alimentation. Cohérent avec le fait que ce LiDAR
fonctionnait auparavant (PrototypeVA, 2025) : connecteur probablement desserré
par les vibrations.

\begin{tcolorbox}[colback=lightgreen,colframe=success,title=Leçon clé]
La rotation de la tête ne prouve que l'\textbf{alimentation}. Une sonde série
brute muette à tous les débits prouve que le défaut est sur la \textbf{liaison de
données} (connecteur/câble), et non dans la configuration logicielle.
\end{tcolorbox}

% =====================================================================
\section{Réparation et validation}
\subsection{Geste de réparation}
Mise hors tension, puis \textbf{ré-emboîtement ferme du connecteur de données}
entre la tête du LiDAR et la carte CP2102 (tirer sur le connecteur, jamais sur
les fils ; respecter le sens), puis rebranchement de l'USB.

\subsection{(\ok) La sonde répond}
\begin{lstlisting}
[115200] GET_INFO   -> a5 5a 14 00 00 00 04 18 1d 01 07 7e e4 ...
[115200] GET_HEALTH -> a5 5a 03 00 00 00 06 00 00 00
\end{lstlisting}
Décodage : descripteur \texttt{A5 5A}, \textbf{modèle 0x18}, firmware
\textbf{1.29}, révision matérielle 7, et statut de santé \texttt{00} = \textbf{OK}.
$\Rightarrow$ le capteur est un \textbf{RPLIDAR A1}, débit \textbf{115200}.

\subsection{(\ok) Le topic \texttt{/scan} est publié}
\begin{lstlisting}[language=bash]
ros2 run rplidar_ros rplidar_composition --ros-args \
  -p serial_port:=/dev/ttyUSB0 -p serial_baudrate:=115200 \
  -p frame_id:=laser_frame -p angle_compensate:=true
ros2 topic hz /scan
\end{lstlisting}
\begin{center}
\begin{tabular}{@{}ll@{}}
\toprule
Paramètre & Valeur \\
\midrule
Santé / mode      & OK / \emph{Sensitivity}, 8 kHz, 12 m, 10 Hz \\
Cadence \texttt{/scan} & $\approx$ 7,4 Hz \\
Points par tour   & $\approx$ 1080 \\
Portée            & 0,15 -- 12 m \\
\texttt{frame\_id} & \texttt{laser\_frame} \\
\bottomrule
\end{tabular}
\end{center}

% =====================================================================
\section{Visualisation dans RViz (côté PC, WiFi)}
\subsection{Disponibilité de RViz}
\begin{lstlisting}[language=bash]
echo $DISPLAY        # -> :1  (session graphique presente)
command -v rviz2     # -> /opt/ros/humble/bin/rviz2  (deja installe)
\end{lstlisting}
RViz\,2 est fourni avec ROS\,2 Humble : aucune installation supplémentaire
n'a été nécessaire, et une session graphique (\texttt{DISPLAY=:1}) était
disponible sur le PC.

\subsection{Maintenir le LiDAR actif : tentatives échouées}
RViz tourne sur le PC ; le LiDAR doit tourner sur la Jetson. Faire survivre un
processus lancé par SSH après la fermeture de la session s'est révélé délicat
sur cette Jetson :
\begin{center}
\begin{tabularx}{\textwidth}{@{}l c X@{}}
\toprule
Méthode & Résultat & Raison \\
\midrule
\texttt{nohup \ldots \& disown}      & \ko & le processus meurt à la fermeture du SSH \\
\texttt{setsid \ldots}               & \ko & idem (SIGHUP) \\
\texttt{tmux}                        & \ko & \texttt{tmux} non installé \\
SSH maintenu ouvert (1\textsuperscript{er} essai) & \ko & auto-kill (voir piège ci-dessous) \\
SSH maintenu ouvert (corrigé)        & \ok & le LiDAR reste publié tant que la session vit \\
\bottomrule
\end{tabularx}
\end{center}

\begin{tcolorbox}[colback=lightorange,colframe=warning,title=Piège : auto-suppression par pkill]
\texttt{pkill -f "rplidar\_composition"} lancé \emph{dans} une commande SSH
correspond à la \textbf{ligne de commande du shell SSH lui-même} (qui contient ce
texte) : il se tue donc tout seul (sortie 255). Solutions : faire le ménage dans
une commande SSH \emph{séparée}, utiliser l'astuce des crochets
(\texttt{pkill -f "[r]plidar"}), ou tuer par PID.
\end{tcolorbox}

\noindent Solution retenue : lancer le LiDAR \emph{au premier plan} dans une
session SSH maintenue ouverte ; tant qu'elle vit, \texttt{/scan} est publié.
\begin{lstlisting}[language=bash]
ssh protova2@10.37.11.28 'bash -lc "
  source /opt/ros/humble/setup.bash
  source ~/protova2_yass/install/setup.bash
  export ROS_DOMAIN_ID=125 RMW_IMPLEMENTATION=rmw_fastrtps_cpp
  unset ZENOH_SESSION_CONFIG_URI ZENOH_ROUTER_CONFIG_URI ZENOH_CONFIG
  exec ros2 run rplidar_ros rplidar_composition --ros-args \
    -p serial_port:=/dev/ttyUSB0 -p serial_baudrate:=115200 \
    -p frame_id:=laser_frame -p angle_compensate:=true"'
\end{lstlisting}

\subsection{Configuration RViz}
Astuce : en fixant \textbf{Fixed Frame = \texttt{laser\_frame}} (le repère du
scan lui-même), aucun arbre de transformations (TF) n'est nécessaire pour
afficher le nuage de points. Un fichier de configuration
(\texttt{scan.rviz}) a été préparé : grille de 1 m, affichage \texttt{LaserScan}
sur \texttt{/scan}, fiabilité \emph{Best Effort} (le LiDAR publie ainsi). Un
script de lancement \texttt{view\_scan\_wifi.sh} positionne l'environnement
(DDS natif, domaine 125) puis ouvre RViz.
\begin{lstlisting}[language=bash]
./view_scan_wifi.sh        # = rviz2 -d scan.rviz, env DDS WiFi
\end{lstlisting}

\subsection{(\ok) Résultat}
RViz affiche, en temps réel, un nuage de points rouges représentant le contour
de la pièce vu de dessus : chaque point est un impact laser, le centre (0,0) est
le LiDAR, chaque carré vaut 1 m. \texttt{LaserScan : Status Ok}, 31 fps.
(« Global Status : Warn » est normal : pas d'arbre TF, sans incidence puisque le
\emph{Fixed Frame} est le repère du scan.)

% =====================================================================
\section{Premier essai de conduite réactive (roues en l'air)}
Une fois le \texttt{/scan} validé, un premier essai de la chaîne de conduite
autonome réactive a été réalisé. \textbf{Précaution :} la voiture était
\emph{surélevée, roues ne touchant pas le sol}, afin d'observer la réaction de la
direction et de la propulsion sans déplacement réel.

\subsection{Ce qui a été lancé}
Tout est démarré par un unique script, \texttt{launch\_autonomous.sh}, exécuté
sur la Jetson. Il lance huit nodes ROS\,2 (DDS natif, domaine 125) :
\begin{table}[H]
\centering
\begin{tabularx}{\textwidth}{@{}c l X l@{}}
\toprule
\# & Node & Rôle & Topic / sortie \\
\midrule
1 & \texttt{rplidar\_composition}        & LiDAR                              & \texttt{/scan} \\
2 & \texttt{detection dist\_min}          & distances mini par secteur        & \texttt{/closest\_object\_front\_distance} \\
3 & \texttt{safety AEB}                   & freinage d'urgence (prio.\ 255)   & \texttt{/cmd\_vel\_aeb} \\
4 & \texttt{navigation follow\_the\_gap}  & conduite réactive (prio.\ 10)     & \texttt{/cmd\_vel\_auto} \\
5 & \texttt{controller g29\_teleop}       & reprise pilote (prio.\ 30)        & \texttt{/cmd\_vel\_G29} \\
6 & \texttt{twist\_mux}                   & arbitrage des priorités           & \texttt{/cmd\_vel} \\
7 & \texttt{controller TX}                & \textbf{envoi au Pico (actionneurs)} & série \texttt{/dev/ttyACM0} \\
8 & \texttt{RX}                           & retour capteurs                   & \texttt{/velocity}, \texttt{/ticks} \\
\bottomrule
\end{tabularx}
\end{table}

\noindent \textbf{Hiérarchie de sécurité} (arbitrée par \texttt{twist\_mux}) :
\[
\text{AEB (255)} \;>\; \text{pilote G29 (30)} \;>\; \text{autonome (10)}
\]
Ainsi l'AEB ou le volant reprennent la main à tout moment sur la conduite
autonome.

\subsection{Comment (maintien du processus par SSH)}
Comme précédemment, le script est lancé \emph{au premier plan} dans une session
SSH maintenue ouverte ; son instruction finale \texttt{wait} garde tous les nodes
vivants tant que la session existe :
\begin{lstlisting}[language=bash]
ssh protova2@10.37.11.28 'bash /home/protova2/launch_autonomous.sh'
\end{lstlisting}

\subsection{(\ok) Résultats observés}
La chaîne fonctionne de bout en bout. Les huit nodes sont actifs et TX a bien
ouvert le port série du Pico :
\begin{lstlisting}
[INFO] [cmd_vel_stamped_to_serial]: Serial ouvert sur /dev/ttyACM0 @ 115200
[INFO] [follow_the_gap]: cap=-82.2  servo=28.0  v=0.095 m/s  avant=1.11 m
[INFO] [cmd_vel_stamped_to_serial]: CTRL,28,0.095   <- trame envoyee au Pico
\end{lstlisting}
\begin{center}
\begin{tabular}{@{}ll@{}}
\toprule
Mesure (vue depuis le PC) & Valeur \\
\midrule
\texttt{/cmd\_vel} (servo, vitesse) & 28 ; $\approx$ 0,09 m/s \\
\texttt{/closest\_object\_front\_distance} & 0,73 m \\
follow\_the\_gap : cap visé & $\approx -82$\textdegree\ (espace libre à droite) \\
TX $\rightarrow$ Pico & \texttt{CTRL,28,0.09} en continu \\
\texttt{/velocity} (encodeur) & $-0{,}0$ \emph{(à investiguer)} \\
\bottomrule
\end{tabular}
\end{center}

\noindent \textbf{Interprétation :} le node vise le côté droit (espace le plus
dégagé dans la pièce), braque le servo à fond (28) et commande une vitesse lente
($\approx$0,09 m/s) — comportement cohérent. Les ordres atteignent
physiquement le Pico. Point à vérifier : \texttt{/velocity} reste à $0$, ce qui
peut signifier que la roue motrice ne démarre pas (vitesse sous le seuil de
l'ESC) ou que l'encodeur ne renvoie rien ; à confirmer par observation directe.

\subsection{(\ok) Arrêt sécurisé}
À la demande, tout a été arrêté proprement, en garantissant l'arrêt des
actionneurs :
\begin{enumerate}[leftmargin=1.6em]
  \item arrêt des nodes par \texttt{pkill} en \textbf{forme crochet}
        (\texttt{pkill -9 -f "[r]plidar"}\ldots), sans aucun nom littéral
        ailleurs dans la commande ;
  \item \textbf{envoi d'une trame neutre au Pico} (direction centrée, vitesse
        nulle) pour ne pas laisser le moteur sur la dernière commande :
\begin{lstlisting}[language=bash]
printf "CTRL,83,0.0\n" > /dev/ttyACM0   # servo=83 (centre), vitesse=0
\end{lstlisting}
\end{enumerate}

\begin{tcolorbox}[colback=lightorange,colframe=warning,title=Deux pièges rencontrés à l'arrêt]
\textbf{(1) Auto-suppression par \texttt{pkill} :} si la commande SSH contient un
nom de node en clair (même ailleurs, p.\,ex.\ dans un \texttt{pgrep} de
vérification), un \texttt{pkill -f} sur ce nom tue le shell SSH lui-même
(sortie 255). $\Rightarrow$ tout mettre en forme crochet, vérifier dans une
commande séparée.\\
\textbf{(2) \texttt{kill -9} et \texttt{trap} :} tuer le script avec \texttt{-9}
n'exécute pas son \texttt{trap} de nettoyage : les nodes enfants survivent
(orphelins). $\Rightarrow$ tuer explicitement chaque node, puis envoyer la trame
neutre au Pico.
\end{tcolorbox}

% =====================================================================
\section{Synthèse}
\begin{enumerate}[leftmargin=1.6em]
  \item Capteur identifié : \textbf{RPLIDAR A1} (CP2102, \texttt{/dev/ttyUSB0},
        115200 baud, firmware 1.29).
  \item Panne diagnostiquée par élimination : \textbf{connecteur de données
        desserré} (tête alimentée mais muette) — \emph{toutes} les causes
        logicielles écartées (débit, droits, ModemManager, sonde brute, DTR/RTS).
  \item Réparation : ré-emboîtement du connecteur $\rightarrow$ sonde renvoie
        \texttt{A5 5A}, santé OK.
  \item \texttt{/scan} publié ($\approx$7,4 Hz, $\approx$1080 pts, 0,15--12 m).
  \item Visualisation RViz sur le PC via WiFi (DDS natif), \emph{Fixed Frame}
        \texttt{laser\_frame}, sans TF.
  \item \textbf{Premier essai de conduite réactive} (roues en l'air) : chaîne
        complète lancée, follow\_the\_gap décide, TX commande le Pico ; arrêt
        sécurisé avec trame neutre. À vérifier : démarrage de la roue motrice
        (\texttt{/velocity}=0).
\end{enumerate}

\paragraph{Livrables.}
\texttt{launch\_autonomous.sh} (chaîne autonome complète),
\texttt{scan.rviz} (config RViz), \texttt{view\_scan\_wifi.sh} (lancement RViz),
\texttt{launch\_lidar.sh} (démarrage LiDAR, débit paramétrable),
\texttt{rplidar\_probe.py} (sonde de diagnostic série).

\paragraph{Prochaine étape.}
Confirmer le démarrage de la propulsion (seuil ESC / encodeur), puis régler les
paramètres de \texttt{follow\_the\_gap} (\texttt{bubble\_radius},
\texttt{gap\_threshold}, \texttt{steer\_gain}, vitesses) ; ensuite seulement,
essai au sol dans un couloir dégagé.

\end{document}
